% Propietario conservado: /Users/ruben/Documents/New project/output/EXTREMA_MEDIA_RAZON_RECIPROCIDAD_ES_EN_20260918/ARTICULO/ES/source/sections/electron.tex
\section{Algebraic structure of the electron and action scale}
\label{x_reciprocidad:sec:electron}

The electronic construction receives the arithmetic state with
trajectory memory produced by APP--TRIT--TPK and its numerical coordinates. At this point in the development,
\(\pi,\varphi,e\) and \(\alpha\) designate the coordinates already obtained from
the same genealogy, in the constructive order set out above. They
are not parameters taken from a physical table. The electronic section is a
subsequent realization of that joint discrete structure of the continuum:
it preserves the cell of origin, the two arithmetic sheets, orientation,
carry, and return memory before producing a scalar.

Three objects should be distinguished from the outset. The first is the central
section and its transport history; the second is the dimensionless
coordinate obtained from it; the third is its realization on an
energy or mass line. Equality between two numerical coordinates does not
identify these three levels. In particular, Euler's number \(e\),
the electronic coordinate \(\varepsilon_e\), and a rest energy
\(E_e\) will have distinct notations.

The aim of this section is to establish the composition
\begin{equation}
 \varepsilon_e^{\beta}
 =\frac{\sqrt3}{4}
  \left[\exp\!\left(\frac{\varphi}{\pi^2}\right)-22\alpha^3\right]
  (1+15\Delta_4)
  R_{\mathrm{act}}^{\,80-54\alpha+6\Delta_4},
 \label{x_reciprocidad:eq:el-formula-completa}
\end{equation}
explaining separately the origin of the prefactor, the two integers of the
basal stage, normalization of \(\Delta_4\), the action ratio and the
return triple. The dimensional chart opens only after
this composition. The result does not consist in making a dimensionless
number acquire units merely through its resemblance to a tabulated
quantity.

The composition brings together factors with different functions. The
following table identifies their meaning and refers to the corresponding
construction; the proofs are developed in the order of their dependencies.

\begin{table}[htbp]
\centering
\small
\setlength{\tabcolsep}{4pt}
\renewcommand{\arraystretch}{1.18}
\begin{tabularx}{\linewidth}{@{}p{0.20\linewidth}Xp{0.24\linewidth}@{}}
\toprule
Factor & Meaning in the composition & Construction and proof\\
\midrule
\(\sqrt3/4\)
 & Norm of the commutator of the projectors on the principal plane.
 & Equation~\eqref{x_reciprocidad:eq:el-prefactor}.\\
\(\exp(\varphi/\pi^2)\)
 & Exponential character of the electronic orientation of the regional quotient.
 & Equation~\eqref{x_reciprocidad:eq:el-bisagra}; subsections~\ref{x_reciprocidad:sec:medida-retorno-hojas}
   and~\ref{x_reciprocidad:sec:operador-retorno-marcado}.\\
\(-22\alpha^3\)
 & Contribution of the regional complement in the central return.
 & Equations~\eqref{x_reciprocidad:eq:el-selector-enteros} and~\eqref{x_reciprocidad:eq:el-censo}.\\
\(1+15\Delta_4\)
 & Torsional normalization composed with regional and tritic multiplicity.
 & Equations~\eqref{x_reciprocidad:eq:el-selector-enteros}, \eqref{x_reciprocidad:eq:el-delta}
   and~\eqref{x_reciprocidad:eq:el-basal}.\\
\(R_{\mathrm{act}}^{\kappa_e}\)
 & Quotient of the action sections, raised to the character determined
   by the two readers of the electronic return.
 & Equations~\eqref{x_reciprocidad:eq:el-ratio}, \eqref{x_reciprocidad:eq:el-triple},
   \eqref{x_reciprocidad:eq:el-kappa} and~\eqref{x_reciprocidad:eq:el-beta}.\\
\bottomrule
\end{tabularx}
\caption{Factors of the electronic composition and provenance of their rules.}
\label{x_reciprocidad:tab:el-factores}
\end{table}


% BEGIN OWNER /Users/ruben/Documents/New project/output/EXTREMA_MEDIA_RAZON_RECIPROCIDAD_ES_EN_20260918/ARTICULO/ES/source/sections/electron_articulacion.tex
\paragraph{Identity of the section and composition of its observables.}
\label{x_reciprocidad:el-ampl-articulacion}
The structural identity of the electron is introduced through the central
section and its transport operations. Inversion of the APP chart
selects the cell $(5,5)$ by its fixed-point property. The two arithmetic
sheets allow the deformations preserving the sum and
product residue to be studied. The result distinguishes the two orientations
$(8,2)$ and $(2,8)$ and preserves the difference of one quotient unit.
This is the first level of information: a cell, an involution, a
preserved fiber and a minimal integer defect. No mass evaluation
is needed to state or prove it.

The passage to arithmetic projections introduces a second level. The
sum and product fibres determine conditional expectations on the
same finite space. Their commutator measures dependence on the
order of these two projections. The norm $\sqrt3/4$ follows from the spectrum
of their pairing; the block attaining this norm also preserves
the quadratic relations proved below. Thus, the scalar
prefactor and the local algebra have a common provenance, but
contain different information. The norm retains the magnitude of the
commutator, while the matrices preserve its orientation and its
composition law.

Areal--volumetric transfer incorporates the regional coordinates
at a third level. The mode calendar determines the incidence
matrix and its quadratic action. The closure mark produces the ratio
$\pi/\varphi^2$; interchanging the two arguments provides the
orientation $\varphi/\pi^2$ used in the electronic exponential
character. The two expressions are related by an explicit
involution. This relationship explains why a regional coordinate may
appear in the electronic equation with a power or in a position
different from that in its first reading: the operation applied
is part of the mathematical datum that is preserved.

The cubic alpha term belongs to the regional complement of a marked
return. Counting its positions yields $27-5=22$, while the
three sheets of each of the five regions yield $3\cdot5=15$.
The proofs preserve the domain of these counts and the
orientation determining their signs. The basal formula therefore
composes an incompatibility norm, a transfer character,
a position deficit, and a torsional memory.
Writing it on a single line expresses the composition of these
operations; explaining it requires reconstructing them in that order.

The transition from the basal coordinate to the complete coordinate preserves
the same central section. The return factor uses the history
records, and the two stated readers recover the triple
$(80,54,6)$ along that trajectory. Multiplication by
$R_{\mathrm{act}}^{80-54\alpha+6\Delta_4}$ modifies the evaluation,
but keeps its origin and antecedents identified. The basal
value constitutes an intermediate stage of the composition; the complete
value also includes this return contribution. The physical
comparison must use the stage specified by the formula being evaluated.

The action scale intervenes subsequently with a different function. A ratio of
two action sections is dimensionless and cancels their common scale.
The absolute realization instead uses an action basis, a
temporal reader, and the map to energy or mass. This separation allows
the construction of the scale and the identity of the electronic section
to be preserved simultaneously. The chronological quantum provides a
scale reference; the central section determines the structural factor
expressed relative to it. The dimensional equations of this
chapter keep that composition explicit.

The electron's relevance to the overall argument arises from this
convergence of operations. The same arithmetic that produced the
regions now determines a central cell, its deformations, and its
projections. The fine-structure constant, constructed earlier, acts
in this composition as a closure coordinate. Return information
determines the complete evaluation, and dimensional realization
allows subsequent comparison. Explaining the electron from its structure
consists in preserving this succession of objects, maps, and values at every
step of the proof.

% END OWNER


\subsection{Center, orientation, and vacancy on a preserved fiber}

The cellular support under consideration is \(\{1,\ldots,9\}^2\), with both
APP evaluations, \(r+c\) and \(rc\). The central inversion is
\begin{equation}
 \rho(r,c)=(10-r,10-c).
 \label{x_reciprocidad:eq:el-inversion-central}
\end{equation}
The center is not chosen by a mass: it is characterized by a property of
the support that precedes its physical reading.

\begin{proposition}[Uniqueness of the center]
The inversion \(\rho\) has a unique fixed point, \((5,5)\).
\end{proposition}
\begin{proof}
The equations \(r=10-r\) and \(c=10-c\) are equivalent to
\(2r=2c=10\). Their unique solution in the nonadic chart is \((5,5)\).
\end{proof}

The preceding property identifies a cell, not a unique deep
history. A route adds the initial orientation choice and its
compatible prolongations. In the chart used here the central oriented
section starts at \((5,5)\) with both positive
orientations; it will be denoted by \(\Gamma_e\). Return of its cellular
projection does not erase the transported orientations or the carry
register of the section with its trajectory memory.

Vacancy is now considered, after the closure of \(\alpha\), as
a local operation on this section. Its arithmetic can be isolated in a
finite lemma; this isolation does not make it a substitute for the
preceding generator of constants.

\begin{definition}[Central fiber and quotient defect]
The central additive fiber is
\[
 \mathcal F_{10}
 =\{(5+t,5-t):t\in\mathbb Z,\ |t|\leq4\}.
\]
In the chart \(n=\rho_9(n)+9q_9(n)\), with
\(\rho_9(n)\in\{1,\ldots,9\}\), the quotient defect relative
to the center is defined as
\[
 v(t)=q_9(25)-q_9(25-t^2),
\]
when both products have the same nonadic residue.
\end{definition}

\begin{lemma}[Minimal central vacancy]
The only points of \(\mathcal F_{10}\) that preserve the
multiplicative residue of the center are
\[
 (5,5),\qquad (8,2),\qquad (2,8).
\]
At the two noncentral points, the quotient defect is exactly one.
\end{lemma}
\begin{proof}
The sum is constantly ten, whereas
\[
 (5+t)(5-t)=25-t^2.
\]
Preservation of the residue is equivalent to \(t^2\equiv0\pmod9\), that is,
to \(3\mid t\). In the integer interval \([-4,4]\), the solutions are
\(t=-3,0,3\). For the two nonzero solutions,
\[
 25=7+9\cdot2,\qquad16=7+9\cdot1.
\]
The residue seven is preserved and one unit of quotient is lost. There is
no other admissible noncentral displacement in that domain, so the
positive defect is minimal. The interchange \(t\mapsto-t\) permutes the
two outputs and preserves the defect.
\end{proof}

Vacancy here is not a numerical zero: it is a difference between two
states with the same residual reading and different quotients. If only
the residue were preserved, that difference would disappear.
Orientation also distinguishes the two output points. Neither the quotient
defect nor that orientation is automatically identified with an electric
charge or a mass; their physical realization requires the corresponding
readers.


% BEGIN OWNER /Users/ruben/Documents/New project/output/EXTREMA_MEDIA_RAZON_RECIPROCIDAD_ES_EN_20260918/ARTICULO/ES/source/sections/centro_electronico_registros.tex
\subsubsection{Central coordinates, orientation, and regional records}
\label{x_reciprocidad:el-centro:registros}

The central structure relates position, evaluation, and transport. At
cell \((5,5)\), the two APP sheets give
\begin{equation}
 \mathcal A(5,5)=\bigl((1,1),(7,2)\bigr),\qquad
 5+5=1+9\cdot1,\qquad5\cdot5=7+9\cdot2.
 \label{x_reciprocidad:el-centro:app}
\end{equation}
The two ones on the additive sheet are, in order, residue and quotient.
Their numerical coincidence preserves two distinct arithmetic functions.
On the sum-\(10\) fiber, the two oriented vacancies satisfy
\begin{equation}
 \mathcal A(8,2)=\mathcal A(2,8)=\bigl((1,1),(7,1)\bigr).
 \label{x_reciprocidad:el-centro:vacancias}
\end{equation}
The complete additive sheet and the multiplicative residue are preserved;
the quotient coordinate of the product decreases by exactly
\(1\). Transport additionally retains which of the two conjugate
positions has been reached.


% BEGIN OWNER /Users/ruben/Documents/New project/output/EXTREMA_MEDIA_RAZON_RECIPROCIDAD_ES_EN_20260918/ARTICULO/ES/source/figures/centro_y_vacancias.tex
\begin{figure}[H]
\centering
\begin{tikzpicture}[x=.67cm,y=.67cm,font=\footnotesize,>=Latex]
 \draw[step=1,black!12,line width=.3pt] (1,1) grid (9,9);
 \draw[->,black!60] (.7,1)--(9.5,1) node[right] {$i$};
 \draw[->,black!60] (1,.7)--(1,9.5) node[above] {$j$};
 \foreach \k in {1,...,9}{
  \node[below=2pt] at (\k,1) {$\k$};
  \node[left=2pt] at (1,\k) {$\k$};
 }
 \draw[black!65] (1,9)--(9,1);
 \draw[<->,line width=.9pt] (2,8)--(8,2);
 \foreach \x/\y in {2/8,5/5,8/2}{
  \filldraw[fill=white,line width=.8pt] (\x,\y) circle (3pt);
 }
 \fill (5,5) circle (1.5pt);
 \node[above right=3pt,fill=white,inner sep=1pt] at (2,8) {$(2,8)$};
 \node[above right=3pt,fill=white,inner sep=1pt] at (5,5) {$(5,5)$};
 \node[above right=3pt,fill=white,inner sep=1pt] at (8,2) {$(8,2)$};
 \node[anchor=west,align=left] at (10.5,7.6)
   {$i+j=10$\\$R^+=1,\quad q^+=1$};
 \node[anchor=west,align=left] at (10.5,5.4)
   {center: $ij=25$\\$R^\times=7,\quad q^\times=2$};
 \node[anchor=west,align=left] at (10.5,3.1)
   {vacancies: $ij=16$\\$R^\times=7,\quad q^\times=1$};
\end{tikzpicture}
\caption{Electronic center and conjugate vacancies in the APP chart.
The diagonal is the additive fiber \(\mathcal F_{10}\); the center
is the fixed point of \((i,j)\mapsto(10-i,10-j)\). The two
shifts of magnitude \(3\) preserve the sum and product
residue and decrease the multiplicative quotient by one unit. The
arrows represent deformations along the fiber.}
\label{x_reciprocidad:el-centro:figura}
\end{figure}

% END OWNER


Central symmetry also characterizes its persistence under chart
changes. If an admissible atlas transformation \(f\) intertwines
inversion, \(f\rho=\rho f\), then
\[
 \rho f(5,5)=f\rho(5,5)=f(5,5).
\]
Uniqueness of the fixed point forces \(f(5,5)=(5,5)\). This is the
internal property making the central location stable: every
transformation of that type preserves the cell before evaluating its
electronic character (Section~\ref{x_reciprocidad:el-centro:registros}).

The central matrix block belongs to evaluation on the two
adjacent positions \(4,5\). Its decomposition, established in the
local development of the kernel, is
\[
 B_c=\begin{pmatrix}7&2\\2&7\end{pmatrix}=9P_++5P_-.
\]
Here \(9\) and \(5\) are eigenvalues of two modes, whereas
\((5,5)\) is a position in the atlas. Concatenations of its
entries give \(72+27=77+22=99\) and
\(77-22=55=54+1\). The first equality preserves the total sum under
two arrangements; the second expresses the defect between diagonal
and antidiagonal readings. These operations precede electronic
composition and preserve their matrix provenance.

The initial orientation is another datum. In the representation of
electronic return, \((h_0,v_0)=(1,1)\) is fixed, that is, positive
advance in both coordinates. Simultaneous inversion every \(27\)
steps, read in windows of \(9\), determines
\[
 \tau_e=(+1,+1,+1,-1,-1,-1,+1,+1,+1,-1,-1,-1).
\]
The segments \((1,1,1)\) belong to this orientation sequence:
each segment comprises three windows in the same direction. The recurrence
in the section on readers proves this record in
\eqref{x_reciprocidad:eq:el-tau}, and its two evaluations recover the triple
\((80,54,6)\) entering
\eqref{x_reciprocidad:eq:el-formula-completa}. The positional coordinate,
APP quotients, and sign sequence are thus connected through
explicit operations.

The signature \(555555\) belongs to a regional reading in another
domain. For a decimal emission \(U=(U_1,\ldots,U_6)\), we write
\(U_j=100d_{1j}+10d_{2j}+d_{3j}\), with three digits, including
leading zeros. The multiplicative-signature reader is
\[
 \mathcal E_\times(U)
   =\bigl([d_{2j}-d_{1j}]_{10}\bigr)_{j=1}^{6}.
\]
In the transversal region
\(U_\pi^\perp=(501,614,498,169,272,272)\), the differences are
\((-5,-5,5,5,5,5)\). Consequently,
\begin{equation}
 \mathcal E_\times(U_\pi^\perp)=(5,5,5,5,5,5),\qquad
 \Psi(U_\pi^\perp)=010211.
 \label{x_reciprocidad:el-centro:firma-transversal}
\end{equation}
The first map preserves a multiplicative signature of the catalog;
the second preserves the ternary word common to the five regions.
Figure \ref{x_reciprocidad:pi-estrella:figura} shows the transversal position of
this emission and its multiplicity \(432\).

\paragraph{Memory required to transport the constant signature.}
The signature \(555555\) additionally provides a verifiable example of
the trajectory information required by TPK. Let
\(\mathcal X_\times=D_9^2\times\{N,E,S,O\}\), with \(324\)
oriented positions in the product channel. Advance \(P\) moves
one square in the stored direction, and the half-turn \(H\)
reverses that direction. In this representation, the reader
\(E_{\rm sig}\) sums the discrete exponents by windows:
\[
 \ell_9(1,2,4,8,7,5)=(0,1,2,3,4,5),\qquad
 \ell_9(3)=\ell_9(6)=\ell_9(9)=0.
\]
During each of the six nine-step windows,
\(\ell_9(\rho_9(ij))\) is evaluated before every product advance, selected
by TRIT at times \(2,5,8\pmod9\). The window sum is
reduced modulo \(10\). After steps \(27\) and \(54\), the
direction is reversed. This rule determines the six components of
\(E_{\rm sig}\) from the initial position and orientation.

Enumeration of this domain produces \(26\) signatures. The fiber of
\(555555\) contains \(24\) oriented positions; after one
advance, their images are distributed among three signatures, with eight
representatives each:
\[
 555177,\qquad555717,\qquad555771.
\]
Three witnesses, in coordinates from \(1\) to \(9\), allow the
distinction to be checked directly:
\[
\begin{array}{c|c|c}
x&E_{\rm sig}(x)&E_{\rm sig}(Px)\\\hline
(3,5,N)&555555&555177\\
(3,4,N)&555555&555717\\
(3,4,S)&555555&555771
\end{array}
\]
Thus, the common value of the signature does not determine its next value:
the transport class within that fiber forms part of the operational
datum. The minimal stable partition is calculated by starting from
signature equivalence \(\sim_0\) and refining through
\[
 x\sim_{n+1}y\quad\Longleftrightarrow\quad
 x\sim_n y,\quad Px\sim_n Py,\quad Hx\sim_n Hy.
\]
The procedure terminates because the domain is finite. Every stable
congruence refining \(\sim_0\) refines every \(\sim_n\), by induction;
its fixed point is therefore the minimal stable congruence required.
The census gives \(26\to28\to28\), with histogram
\(27\cdot8+108=324\). The only signature that splits is
\(555555\), into three classes of eight elements. The accompanying certificate
reproduces the enumeration, the three classes, and stability under both
operators (Section~\ref{x_reciprocidad:el-centro:registros}).

The multiplicity \(24\) here belongs to the product-channel fiber;
the multiplicity \(432\) belongs to the complete regional emission
\(U_\pi^\perp\). The first proves why the signature needs transport
memory; the second places that signature within the joint catalog of
the two sheets. Both readings participate in the same architecture,
with specified domains and operations.

This distinction of domains makes it possible to bring together the occurrences of unity
and the center with their respective meanings:
\begin{center}
\small
\begin{tabularx}{\textwidth}{@{}lX@{}}
\toprule
Registro & Object and function \\
\midrule
\((5,5)\) & Fixed cell of the central inversion of the APP atlas.\\
\((R^+,q^+)=(1,1)\) & Residue and quotient of the central additive evaluation.\\
\((h_0,v_0)=(1,1)\) & Initial orientation of the two shifts.\\
\((+1,+1,+1)\) & Three consecutive windows of \(\tau_e\) with positive orientation.\\
\(c=111111\) & Datum of the fifth nonadic boundary, Section \ref{x_reciprocidad:mono-recta:desarrollo}.\\
\(\mathcal E_\times=555555\) & Signature of the transversal region of \(\pi\), with six positions.\\
\bottomrule
\end{tabularx}
\end{center}

Electronic coordination is obtained by composing these levels in their
dependency order. APP determines the center and evaluations;
TRIT selects regimes and orientations; TPK prolongs the trajectory
and preserves its records. The previously generated regions provide
the characters \(\pi,\varphi,e\); the dodecaphase record provides
\(\alpha\); the electronic readers determine the factors of
equation \eqref{x_reciprocidad:eq:el-formula-completa}. This composition preserves
the identity of the section throughout its different
representations, up to dimensional evaluation and subsequent
metrological comparison.

% END OWNER


\subsection{Commutator norm of the arithmetic projections}

The electronic prefactor is obtained without using \(\alpha\), torsion,
or a dimensional unit. The domain is the third realization of APP,
\(\mathcal A_3=\{1,\ldots,9\}^3\), equipped with the uniform measure.
We define
\[
 \Sigma(x)=\rho_9(x_1+x_2+x_3),\qquad
 \Pi(x)=\rho_9(x_1x_2x_3).
\]
The conditional expectations \(P_\Sigma\) and \(P_\Pi\) are the
orthogonal projections onto functions constant on each additive
and multiplicative fiber, respectively. Thus, for a function \(f\),
\[
 (P_\Sigma f)(x)
 =\frac{1}{|\Sigma^{-1}(\Sigma(x))|}
   \sum_{y:\,\Sigma(y)=\Sigma(x)}f(y),
\]
and analogously for \(P_\Pi\). The two projections act on the
same space; they do not represent two independent arithmetic supports.

\begin{lemma}[Two projections and principal angles]
If \(s\in[0,1]\) is a singular value of the pairing between orthonormal
bases of the ranges of two projections \(P,Q\), its nontrivial
block contributes to \(\|[P,Q]\|\) with
\(s\sqrt{1-s^2}\).
\end{lemma}
\begin{proof}
On the corresponding principal plane, one can write
\[
 P=\begin{pmatrix}1&0\\0&0\end{pmatrix},\qquad
 Q=\begin{pmatrix}
 s^2&s\sqrt{1-s^2}\\
 s\sqrt{1-s^2}&1-s^2
 \end{pmatrix}.
\]
Their commutator is
\[
 [P,Q]=s\sqrt{1-s^2}
 \begin{pmatrix}0&1\\-1&0\end{pmatrix}.
\]
The final matrix is orthogonal, so its norm is one. Common
or orthogonal blocks have zero commutator. The orthogonal decomposition
into principal planes gives the claim and reduces the total norm to the maximum
of these contributions.
\end{proof}

\begin{proposition}[Exact incompatibility norm]
On \(\ell^2(\mathcal A_3)\),
\begin{equation}
 \|[P_\Sigma,P_\Pi]\|=\frac{\sqrt3}{4}.
 \label{x_reciprocidad:eq:el-prefactor}
\end{equation}
\end{proposition}
\begin{proof}
Let \(N_{ab}=|\Sigma^{-1}(a)\cap\Pi^{-1}(b)|\). Integer enumeration
of the three symbols yields
\[
 N=9\begin{pmatrix}
 0&1&1&0&1&1&0&1&4\\
 1&0&1&1&0&1&1&0&4\\
 1&0&2&0&1&2&0&0&3\\
 0&1&1&0&1&1&0&1&4\\
 1&0&1&1&0&1&1&0&4\\
 0&0&2&1&0&2&0&1&3\\
 0&1&1&0&1&1&0&1&4\\
 1&0&1&1&0&1&1&0&4\\
 0&1&2&0&0&2&1&0&3
 \end{pmatrix}.
\]
Each row sums to \(81\); the column sums are
\[
 (m_1,\ldots,m_9)=(36,36,108,36,36,108,36,36,297).
\]
The normalized indicator functions of the fibers have pairing
matrix \(C_{ab}=N_{ab}/\sqrt{81m_b}\). Consequently,
\[
 M=CC^{\mathsf T},\qquad
 M_{ac}=\sum_{b=1}^9\frac{N_{ab}N_{cb}}{81m_b}
\]
is a rational matrix whose spectrum consists of the squares of the
required singular values.

Exact multiplication shows that \((1,\ldots,1)\),
\((1,-1,0)^3\), and \((1,1,-2)^3\) are eigenvectors with eigenvalues
\(1\), \(1/4\), and \(7/132\), respectively. The subspace of
vectors supported on \(\{3,6,9\}\) whose sum is zero has dimension
two and eigenvalue \(1/18\). The zero-sum vectors supported on
\(\{1,4,7\}\) or \(\{2,5,8\}\) form a four-dimensional
kernel. These subspaces are independent and their dimensions sum to
nine. The complete spectrum is therefore established:
\[
 \operatorname{spec}(M)=
 \left\{1,\frac14,\frac7{132},\frac1{18},\frac1{18},0,0,0,0\right\}.
\]
The maximum of \(\lambda(1-\lambda)\) on this set is \(3/16\),
attained at \(\lambda=1/4\). The preceding lemma gives
\(\sqrt{3/16}=\sqrt3/4\).
\end{proof}

The mode \((1,-1,0)^3\) is precisely the value distribution of the
tritic phase selector in the nonadic chart. The scalar norm retains
a measure of incompatibility, but does not by itself preserve the
orientation of that mode or reconstruct the two projectors. That memory
remains in the route state accompanying the scalar reading.

\begin{proposition}[Algebra of the local electronic block]
On the principal plane attaining \eqref{x_reciprocidad:eq:el-prefactor}, let
\[
 P=\begin{pmatrix}1&0\\0&0\end{pmatrix},\qquad
 Q=\begin{pmatrix}1/4&\sqrt3/4\\\sqrt3/4&3/4\end{pmatrix},\qquad
 S=2P-I,\qquad J=\frac4{\sqrt3}[P,Q].
\]
Then \(S^2=I\), \(J^2=-I\), \(SJ=-JS\), and the real algebra
generated by \(S,J\) is \(M_2(\mathbb R)\), a realization of
\(\mathrm{Cl}_{1,1}\). Moreover,
\[
 n_\pm=\frac{S\pm J}{2},\qquad
 n_\pm^2=0,\qquad n_+n_-+n_-n_+=I.
\]
\end{proposition}
\begin{proof}
We have \(S=\operatorname{diag}(1,-1)\) and
\(J=\bigl(\begin{smallmatrix}0&1\\-1&0\end{smallmatrix}\bigr)\).
Multiplication verifies the three relations. The matrices
\(I,S,J,SJ\) are linearly independent and form a basis of
\(M_2(\mathbb R)\), proving the claim about the algebra.
Finally, \((S\pm J)^2=S^2+J^2\pm(SJ+JS)=0\), and the sum of the
cross products is \((S^2-J^2)/2=I\).
\end{proof}

In this block, an elliptic direction, a hyperbolic direction, and two
parabolic null directions coexist. The result describes a local
matrix representation of the arithmetic of the two sheets. It does not identify all of APP
with \(M_2(\mathbb R)\), or with a quantum Latin square; such
claims would require other domains and other relations. Nor does the
exponential belong exclusively to the parabolic regime: it can
be applied to any of the generators of this block.


% BEGIN OWNER /Users/ruben/Documents/New project/output/EXTREMA_MEDIA_RAZON_RECIPROCIDAD_ES_EN_20260918/ARTICULO/ES/source/sections/helicidad_electron.tex
% Adición focal propuesta; no sustituye el álgebra electrónica existente.
% Inserción: después de la proposición «Álgebra del bloque electrónico local»
% y su párrafo explicativo, antes de «Transferencia entre hojas...».
% Procedencia y alcance: LOCALIZADORES_Y_ALCANCE.md, en esta carpeta.
\subsection{Spinor representation of the electronic block and helicity}
\label{x_reciprocidad:sec:el-representacion-espinorial}

The electronic structure has a representation preserving more
information than its scalar evaluation. The principal plane of the two
three-dimensional APP projectors, selected by the tritic mode
\((1,-1,0)^3\), has already provided the operators \(S\) and \(J\).
The former is a symmetric involution; the latter is the normalized
commutator and preserves sheet orientation. Their composition makes it possible
to construct explicitly a rotation representation and separate its
two helicity components. This construction acts on the electronic
block; the TPK route record and its memory remain as data
of the block's provenance.

Let \(E_e\) be the real principal plane, equipped with the inner product inherited
from the conditional expectations, and let
\(\mathscr S_e=E_e\otimes_{\mathbb R}\mathbb C\) be its complexification.
The operators already obtained satisfy
\[
 S^*=S,\qquad J^*=-J,\qquad
 S^2=I,\qquad J^2=-I,\qquad SJ=-JS.
\]
Complexification extends the representation field. It preserves the
real operators and allows the scalar unit \(i\) to be distinguished from the
endomorphism \(J\), although both square to \(-1\) in their
respective domains.

\begin{proposition}[Hermitian triple induced by the two sheets]
\label{x_reciprocidad:prop:el-terna-hermitica}
The operators
\begin{equation}
 \sigma_1=SJ,\qquad \sigma_2=-iJ,\qquad \sigma_3=S
 \label{x_reciprocidad:eq:el-terna-hermitica}
\end{equation}
are Hermitian and traceless. They satisfy
\begin{equation}
 \sigma_a\sigma_b
 =\delta_{ab}I+i\sum_{c=1}^3\epsilon_{abc}\sigma_c,
 \qquad
 \frac12\operatorname{tr}(\sigma_a\sigma_b)=\delta_{ab},
 \label{x_reciprocidad:eq:el-algebra-terna}
\end{equation}
where \(\epsilon_{123}=1\). They therefore form an orthonormal basis
of the real space \(V_e\) of traceless Hermitian endomorphisms of
\(\mathscr S_e\), for the inner product
\(\langle X,Y\rangle=\frac12\operatorname{tr}(XY)\).
\end{proposition}
\begin{proof}
The preceding relations give \((SJ)^*=(-J)S=SJ\) and
\((-iJ)^*=-iJ\). The squares of all three matrices are \(I\).
Furthermore,
\[
 (SJ)(-iJ)=iS,\qquad (-iJ)S=iSJ,\qquad S(SJ)=J=i(-iJ).
\]
Reversing the order changes the sign. This proves the first identity
in \eqref{x_reciprocidad:eq:el-algebra-terna}. In the orthonormal basis of the principal plane,
the matrices are
\[
 \sigma_1=\begin{pmatrix}0&1\\1&0\end{pmatrix},\quad
 \sigma_2=\begin{pmatrix}0&-i\\i&0\end{pmatrix},\quad
 \sigma_3=\begin{pmatrix}1&0\\0&-1\end{pmatrix}.
\]
Their trace is zero, and taking traces in the product identity proves
orthonormality. Every traceless Hermitian endomorphism has the form
\(\bigl(\begin{smallmatrix}z&x-iy\\x+iy&-z\end{smallmatrix}\bigr)\),
with real \(x,y,z\), and thus belongs to their real span.
\end{proof}

In the usual matrix representation, this triple is called the
Pauli matrices. Their use here follows the construction of
\(P,Q,S,J\): they are obtained through the products
\eqref{x_reciprocidad:eq:el-terna-hermitica}. The direction space of this
representation is the unit sphere of \(V_e\), whose positive form
has just been constructed through the trace.

\begin{proposition}[Rotation generators and directional decomposition]
\label{x_reciprocidad:prop:el-helicidad}
Define \(L_a=\sigma_a/2\). Then
\begin{equation}
 [L_a,L_b]=i\sum_c\epsilon_{abc}L_c,
 \qquad \sum_{a=1}^3L_a^2=\frac34 I.
 \label{x_reciprocidad:eq:el-generadores-spin}
\end{equation}
For \(n=(n_1,n_2,n_3)\) with \(\sum_an_a^2=1\), let
\[
 H_e(n)=\sum_an_a\sigma_a,\qquad
 h_e(n)=\frac12H_e(n),\qquad
 \Pi_\pm(n)=\frac12\bigl(I\pm H_e(n)\bigr).
\]
The linear formula \(H_e(v)=\sum_av_a\sigma_a\) is also used
for \(v\) of arbitrary norm.
One has
\begin{equation}
 \mathscr S_e=\operatorname{im}\Pi_+(n)
             \mathbin{\oplus}\operatorname{im}\Pi_-(n),
 \qquad
 h_e(n)\Pi_\pm(n)=\pm\frac12\Pi_\pm(n).
 \label{x_reciprocidad:eq:el-descomposicion-helicidad}
\end{equation}
The two summands are orthogonal complex lines. Rotation
transport around \(n\) is
\begin{equation}
 U_n(\theta)=\exp\!\left(-\frac{i\theta}{2}H_e(n)\right)
 =\cos\frac\theta2\,I-i\sin\frac\theta2\,H_e(n),
 \qquad
 U_n(2\pi)=-I,\quad U_n(4\pi)=I.
 \label{x_reciprocidad:eq:el-retorno-spin}
\end{equation}
\end{proposition}
\begin{proof}
The product identity of the triple gives the commutators in
\eqref{x_reciprocidad:eq:el-generadores-spin}; their squares sum to \(3I/4\).
Antisymmetric terms cancel when multiplying
\(H_e(n)\) by itself, so that \(H_e(n)^2=I\).
It follows that
\[
 \Pi_\pm(n)^2=\Pi_\pm(n),\qquad
 \Pi_+(n)\Pi_-(n)=0,\qquad \Pi_+(n)+\Pi_-(n)=I.
\]
They are Hermitian and have trace \(1\); their rank is therefore \(1\).
Multiplication by \(H_e(n)/2\) proves
\eqref{x_reciprocidad:eq:el-descomposicion-helicidad}. The exponential series
splits into even and odd powers because \(H_e(n)^2=I\), producing
\eqref{x_reciprocidad:eq:el-retorno-spin}. It also directly proves
unitarity and the law \(U_n(\theta+\psi)=U_n(\theta)U_n(\psi)\).

Angular normalization is verified on \(V_e\). For every
\(v\in\mathbb R^3\), the same product identity gives
\[
 U_n(\theta)H_e(v)U_n(\theta)^*
 =H_e\!\left(v\cos\theta+(n\mathbin{\times}v)\sin\theta
       +n(n\mathbin{\cdot}v)(1-\cos\theta)\right).
\]
The induced action is a rotation of angle \(\theta\) and axis \(n\).
The factor \(1/2\) in the spinor generator is therefore the one
corresponding to that angular parameterization. A complete rotation restores
the direction in \(V_e\) and changes the sign of the vector in
\(\mathscr S_e\); two complete rotations restore both.
\end{proof}

The representation is irreducible: a complex subspace invariant under
the three matrices would be invariant under all of \(M_2(\mathbb C)\), the
algebra they generate. Its only invariant subspaces are \(0\) and
\(\mathscr S_e\). The relations
\eqref{x_reciprocidad:eq:el-generadores-spin} thus identify the spin-\(1/2\)
representation. If \(n\) is the propagation direction in a
kinematic realization, \(h_e(n)\) is its dimensionless helicity, with
values \(\pm1/2\). The kinematic direction is declared when that
realization is made; the two projectors are defined for every direction
before a particular propagation state is selected.

\paragraph{Preservation of the electronic structure and types of inversion.}
Direction inversion has an exact law:
\[
 H_e(-n)=-H_e(n),\qquad \Pi_\pm(-n)=\Pi_\mp(n).
\]
This operation exchanges helicity labels relative to a
direction. Charge conjugation, by contrast, acts on the oriented
signature of the material route; cellular inversion acts on the
atlas coordinates. Their domains and actions remain distinct.
Uniqueness of cell \((5,5)\) is compatible with the two lines in
\eqref{x_reciprocidad:eq:el-descomposicion-helicidad}: an atlas cell is not a
vector in the spinor representation it carries.

The electronic evaluation \(\varepsilon_e^\beta\) preserves its
formula and readers. On \(\mathscr S_e\), its scalar realization
is \(\varepsilon_e^\beta I\), which commutes with the generators and with
\(\Pi_\pm(n)\). Helicity decomposition adds rotational
information without modifying the evaluation or turning spin into a
multiplicative mass coefficient.

Chirality corresponds to the grading of the relativistic representation
used: an involution \(\Gamma_{\rm ch}\) determines its
projectors \((I\mp\Gamma_{\rm ch})/2\). The helicity in
\eqref{x_reciprocidad:eq:el-descomposicion-helicidad} additionally depends on \(n\).
For its part, the holonomy of a real Möbius band describes the change
of sign of a line when traversing a loop in its base. The law of that line,
the spinor return \eqref{x_reciprocidad:eq:el-retorno-spin}, and nonadic phase holonomy
with memory advance are transport laws on specified
objects. Their comparison preserves those objects: the helical form
of a trajectory does not replace the representation just
constructed or by itself fix chirality or charge.

% END OWNER


\subsection{Transfer between sheets and coefficients of the basal stage}

The primitive tritic calendar visits the modes
\((\Sigma,\Pi,\Pi)\). Counting consecutive pairs, including the
closure of the block, yields one transition \(\Sigma\to\Pi\), one
\(\Pi\to\Pi\), and one \(\Pi\to\Sigma\). Their primitive incidence is
\[
 F_{\mathrm{av}}=\begin{pmatrix}0&1\\1&1\end{pmatrix}.
\]
This matrix realization is obtained after the calendar. It is not
used to replace the coinductive reader that has already generated \(\varphi\).
Its characteristic polynomial is \(X^2-X-1\), and recognition of the
two roots as \(\varphi\) and \(-\varphi^{-1}\) describes the action of
that output in modal transfer. The quadratic action has eigenvalues
\(\varphi^2,-1,\varphi^{-2}\); the last is the positive stable mode.
The regional mark \(\pi\), applied once, gives
\(\pi/\varphi^2\).

The trajectory space, its measure, and the closure operator are
constructed in subsections~\ref{x_reciprocidad:sec:medida-retorno-hojas}
and~\ref{x_reciprocidad:sec:operador-retorno-marcado}. Formula
\eqref{x_reciprocidad:eq:retorno-marcado-finito} preserves its provenance as a quotient
of returns; the electronic application subsequently composes this reading
with the following involution. The regional signature \(555555\) accompanies
the closure region, whose transport preserves the state with memory;
the centre \((5,5)\) fixes the cell of origin of the oriented electronic
section. Both kinds of information converge in the composition and retain
their respective domains.

\begin{lemma}[Interchange of the two orientations]
Let \(\mathcal R(x,y)=x/y^2\) and \(\mathscr S_{\pi\varphi}(x,y)=(y,x)\), for \(x,y>0\).
Then
\begin{equation}
 \begin{aligned}
 \mathscr S_{\pi\varphi}^{\,2}&=I,\qquad
 \mathcal R(\pi,\varphi)=\frac{\pi}{\varphi^2},\\
 (\mathcal R\circ\mathscr S_{\pi\varphi})(\pi,\varphi)
 &=\frac{\varphi}{\pi^2}
 =\frac1{\varphi^3(\pi/\varphi^2)^2}.
 \end{aligned}
 \label{x_reciprocidad:eq:el-bisagra}
\end{equation}
\end{lemma}
\begin{proof}
Applying the interchange twice gives the identity. Moreover,
\(\varphi^3(\pi/\varphi^2)^2=\pi^2/\varphi\); taking its inverse proves
the last equality.
\end{proof}

The electronic reader uses the orientation \(\varphi/\pi^2\),
followed by the exponential character. The preceding identity explains the
interchange, but does not replace the rule that chooses that character or
by itself produce the integers of the correction. These have a different provenance.

The phase-compatible central return has minimal length
twenty-seven. With the phase and sheet record, its finite space is
\[
 \mathcal K_e=\mathbb Z/9\mathbb Z\times\mathbb F_3.
\]
The five regions of \(\pi\) are distinct genealogical identities,
not five interchangeable copies of a visible expansion. In the regional
gauge fixed by the construction, their injection into the return occupies
the positions
\begin{equation}
 (4,0),\quad(5,0),\quad(6,0),\quad(6,1),\quad(6,2).
 \label{x_reciprocidad:eq:el-cinco-lugares}
\end{equation}
The transverse orientation, mutated return, and three positive
regions remain separated by that record. Admissible changes
of gauge transport the labels jointly; they do not reselect
the positions from an electronic value.

\begin{proposition}[Deficit and retained memory]
Given the central return and the preceding regional injection, the oriented
boundary deficit and torsional memory are
\begin{equation}
 n_e=-|\mathcal K_e\setminus\iota(\mathscr R_\pi)|=-22,
 \qquad
 m_e^{\mathrm{tor}}=|\mathscr R_\pi\times\mathbb F_3|=15.
 \label{x_reciprocidad:eq:el-selector-enteros}
\end{equation}
\end{proposition}
\begin{proof}
The five positions in \eqref{x_reciprocidad:eq:el-cinco-lugares} are distinct.
Thus, \(|\mathcal K_e|=9\cdot3=27\), and the complement of the image
has cardinality \(27-5=22\). The orientation convention records
the deficit removed from the boundary as negative. The retained memory has
three sheets per region, so its positive cardinality is
\(5\cdot3=15\).
\end{proof}

The same pair has the integer control of the family's APP chart:
\begin{equation}
 396=18+72+108+198,
 \qquad \frac{396}{18}=22,
 \qquad \frac{270}{18}=15=3\cdot5=\binom62.
 \label{x_reciprocidad:eq:el-censo}
\end{equation}
The weight \(18\) belongs to the central \(2\times2\) kernel, and \(270\)
to the declared global channel. These equalities check the normalization
of the count; the return argument also fixes the signs. The
quotient of the alternating rings,
\(396/(18+108)=22/7\), is rational and is not identified with \(\pi\).
Nor does \(\binom62\) mean \(6/2\): it counts the unordered pairs
of a six-element set.

Evaluation of the deficit in the third symmetric realization uses
\(\alpha^3\). This power is the product of three copies of the
closure scalar already obtained; it does not assert a universal rule
associating any dimension with the same power of \(\alpha\).
The specified domain and reader are essential to the composition.

\subsection{Global torsion and the basal coordinate}

\begin{definition}[Electronic torsional normalization]
The reduced defect \(d_4\) and the complete reader
\(\Delta_4\) are distinguished:
\begin{equation}
 d_4=\frac{\pi-e\log\pi}{270},\qquad
 \Delta_4=d_4\left(1+\frac{\pi}{729}\right).
 \label{x_reciprocidad:eq:el-delta}
\end{equation}
The denominator \(270\) corresponds to the census channel, whereas the
factor \(1+\pi/729\) retains the coupling with the cubic cell of
\(729\) positions. The electronic reading uses \(\Delta_4\) both in
its basal stage and in the return exponent.
\end{definition}

This definition acts on previously generated \(\pi\) and \(e\). The
logarithm is a subsequent operation in their real realization, not a
retrospective input of conventional values. Torsion is global:
it is not determined independently for each section by an observed
mass.

\begin{lemma}[The two normalizations are not interchangeable]
One has
\begin{equation}
 \Delta_4-d_4=\frac{\pi}{729}d_4.
 \label{x_reciprocidad:eq:el-delta-diferencia}
\end{equation}
In particular, \(d_4>0\) and \(\Delta_4>d_4\) for the coordinates
under consideration.
\end{lemma}
\begin{proof}
The identity follows by distributivity. For positivity, the function
\(f(x)=x-e\log x\), on \(x>0\), has derivative \(1-e/x\) and
second derivative \(e/x^2>0\). Its unique minimum is at \(x=e\), where
it is zero. Since \(\pi\ne e\), we have \(f(\pi)>0\). Both factors
in \eqref{x_reciprocidad:eq:el-delta} are positive and \(1+\pi/729>1\).
\end{proof}

\begin{remark}[Preservation of the normalization]
Removing the second factor in \eqref{x_reciprocidad:eq:el-delta} changes the reader; it
is not an algebraic simplification. If \(d_4\) is introduced as an
auxiliary coordinate, the same object must be written as
\(\Delta_4=(1+\pi/729)d_4\) in every occurrence. Keeping the
integers \(15\) and \(6\) and replacing only \(\Delta_4\) by \(d_4\)
produces a different formula.
\end{remark}

\begin{definition}[Basal electronic coordinate]
The composition of incompatibility, modal transfer, deficit, and retained memory
is
\begin{equation}
 \varepsilon_e^{(0)}
 =\frac{\sqrt3}{4}
  \left[\exp\!\left(\frac{\varphi}{\pi^2}\right)-22\alpha^3\right]
  (1+15\Delta_4).
 \label{x_reciprocidad:eq:el-basal}
\end{equation}
\end{definition}

The equation combines three operations of different kinds. The norm
\(\sqrt3/4\) quantifies the incompatibility of the sheet
projections; its origin is a principal angle of the APP fibers.
The exponential acts on the orientation \(\varphi/\pi^2\) of the
transfer reader. The cubic subtraction records the \(22\) complementary
positions of the marked return, and the torsional factor preserves the
\(15\) components corresponding to three sheets per region. The
composition preserves these domains even though its final evaluation is scalar.

The presence of \(\pi\), \(\varphi\), and \(e\) thus has operational
meanings. Euler's number enters through the exponential and
the logarithmic normalization of \(\Delta_4\); \(\varphi\) and \(\pi\)
enter through oriented transfer. The constant \(\alpha\)
acts after its own construction through the cubic closure
term. The chronological return must still act on this basal
coordinate. Separating the stages makes it possible to explain what changes when
an orientation, normalization, or record is modified, while preserving the central
structure that identifies the electronic section.

The parentheses in this equation are part of its content. The term
\(-22\alpha^3\) is subtracted \emph{outside} the exponential; the memory
\(1+15\Delta_4\) multiplies the entire bracket. Placing the deficit in
the exponent or applying a different torsion to it would alter the composition.

\begin{proposition}[Basal positivity]
For \(0<\alpha<10^{-2}\), \(\varphi>0\), and the normalization
\eqref{x_reciprocidad:eq:el-delta}, we have \(\varepsilon_e^{(0)}>0\).
\end{proposition}
\begin{proof}
The exponential is greater than one, and
\(22\alpha^3<22\cdot10^{-6}<1\). The bracket is therefore positive.
Also, \(1+15\Delta_4>1\) and \(\sqrt3/4>0\).
\end{proof}

The subsequent evaluation of \eqref{x_reciprocidad:eq:el-basal} begins with
\[
 \varepsilon_e^{(0)}=0.5109989224880660725\ldots.
\]
This number is still a dimensionless coordinate and is not the
final stage: the return memory of the route itself remains to be applied.

\subsection{Determinantal action scale and decimal valuation}

The absolute scale and relative correction arise from distinct objects. The former uses the local block of the step-three dodecaphase register,
\begin{equation}
 H_4=\begin{pmatrix}
 1&1&1&1\\1&1&-1&-1\\1&-1&1&-1\\1&-1&-1&1
 \end{pmatrix},\qquad
 G_H=\mathbb Z^4/H_4\mathbb Z^4.
 \label{x_reciprocidad:eq:el-h4}
\end{equation}
The operation uses a local quotient. The three blocks of the complete register are not replaced by tensorization that would multiply this reader's number of sheets again.

\begin{lemma}[Number of sheets of the local quotient]
The Smith normal form of \(H_4\) is
\(\operatorname{diag}(1,2,2,4)\). Consequently,
\[
 G_H\simeq\mathbb Z/2\mathbb Z\oplus\mathbb Z/2\mathbb Z
 \oplus\mathbb Z/4\mathbb Z,\qquad |G_H|=16.
\]
\end{lemma}
\begin{proof}
The greatest common divisor of the entries is one. Second-order minors are even and one has value \(-2\), so their greatest common divisor is two. Since \(H_4H_4^{\mathsf T}=4I\) and \(\det H_4=-16\), the cofactors are \(\pm4\); the third determinantal divisor is four. The fourth is \(16\). Successive quotients of the divisors \(1,2,4,16\) are \(1,2,2,4\). Their product gives the quotient's order.
\end{proof}

\begin{definition}[Determinantal action reader]
The multiplicative norm of the scalar section \(\alpha\), constant on the sheets of \(G_H\), and one application of self-scaling \(\varphi\) produce
\begin{equation}
 \Sigma_H=\varphi\,\mathrm N_{G_H}(\alpha),\qquad
 \mathrm N_{G_H}(\alpha)=\prod_{g\in G_H}\alpha=\alpha^{16}.
 \label{x_reciprocidad:eq:el-lector-determinantal}
\end{equation}
The associated decimal chart determines \(\operatorname{ord}_{10}(\Sigma_H)=\lfloor\log_{10}\Sigma_H\rfloor\).
\end{definition}

The sixteenth power thus comes from the quotient's number of sheets, once this local reader is fixed. The factor \(\varphi\) acts once on the basis of the action line. Smith calculation alone would not select any possible functional on the quotient: the multiplicative norm and self-scaling are part of the declared operation.

\begin{theorem}[Decade of the action section]
The reader \eqref{x_reciprocidad:eq:el-lector-determinantal}, applied to HMT coordinates already generated, satisfies
\begin{equation}
 10^{-34}<\Sigma_H<10^{-33},\qquad
 \operatorname{ord}_{10}(\Sigma_H)=-34.
 \label{x_reciprocidad:eq:el-decada}
\end{equation}
\end{theorem}
\begin{proof}
The rational intervals, much coarser than the available cylinders, suffice:
\[
 \frac{1618}{1000}<\varphi<\frac{1619}{1000},\qquad
 \frac{7297}{10^6}<\alpha<\frac{7298}{10^6}.
\]
The function \((x,a)\mapsto xa^{16}\) is increasing in both positive
arguments. The two integer inequalities
\[
 1618\cdot7297^{16}>10^{65},\qquad
 1619\cdot7298^{16}<10^{66}
\]
imply
\[
 10^{-34}
 <\frac{1618}{1000}\left(\frac{7297}{10^6}\right)^{16}
 <\Sigma_H
 <\frac{1619}{1000}\left(\frac{7298}{10^6}\right)^{16}
 <10^{-33}.
\]
The definition of decimal order completes the proof. Neither an SI unit nor a reference mass enters these inequalities.
\end{proof}

The normalized section begins with \(\Sigma_H=1.0462438381\ldots\times10^{-34}\). Its mantissa is not identified with the action coordinate before or after return. The determinantal reader fixes decimal order; the two return coordinates are obtained through the following readers.


\subsection{Action sections and dimensionless return quotient}

Let \(\mathcal L_S\) be an oriented action line and \(\mathcal U_S\)
its positive basis. The electronic composition uses two sections of
this line and their dimensionless quotient. Its construction requires
the determinantal decade already obtained, the fifth-order action character,
and the regional contribution of the return.

\subsubsection{Action character and provenance of its coefficients}
\label{x_reciprocidad:sec:revision-planck}

The reduced Planck constant, \(\hbar\), is the physical object motivating realization of the angular-action section. Its position in the electronic dependency must be specified at two levels. The determinantal norm \(\Sigma_H\) determines the scale's decade. The character \(H_5\) and regional return then determine the action-section coordinates. This separation allows the result needed for the electron to be presented without anticipating complete developments of other Dimensional Mechanics sectors.

The provenance of the fifth-order character can be made explicit through invariants already fixed on the support. In the central chart, consider
\[
 B_c=\begin{pmatrix}7&2\\2&7\end{pmatrix},
 \qquad \lambda_+=9,\quad\lambda_-=5.
\]
The calendar has length \(12\), the orbit of step \(3\) has length \(4\), and the relevant census quotient is \(104976/1944=54\). The coefficients of the character satisfy
\begin{equation}
 \frac9{16}=\frac{\lambda_+}{4^2},\quad
 \frac59=\frac{\lambda_-}{\lambda_+},\quad
 \frac7{48}=\frac{\operatorname{tr}(B_c)/2}{4\cdot12},\quad
 \frac1{54}=\frac{1944}{104976}.
 \label{x_reciprocidad:eq:revision-h5-coeficientes}
\end{equation}
Assignment of these invariants to degrees \(2,3,4,5\), with the alternating signs of \eqref{x_reciprocidad:eq:el-h5}, belongs to the definition of the analytic character used. Making this explicit is essential: the set of cardinalities, considered without the character and the ordering of its degrees, would admit other expressions. The constructive content lies in composition of the support, orientation and declared analytic rule.

On this basis, the fifth-order character and regional return term
have the expressions
\begin{align}
 H_5={}&\frac12\sqrt{\frac{1000\alpha}{\varphi}}-\alpha
   +\frac9{16}\alpha^2-\frac59\alpha^3
   +\frac7{48}\alpha^4-\frac1{54}\alpha^5,
   \label{x_reciprocidad:eq:el-h5}\\
 D_A={}&\exp\!\left(-\frac{100\pi\alpha}{9}\right),
   \label{x_reciprocidad:eq:el-da}\\
 \mathcal C_\pi(\alpha)={}&169\alpha^6+
       \frac{D_A\alpha^7}{1-D_A\alpha},
   \label{x_reciprocidad:eq:el-cpi}\\
 \eta_{\mathrm{ret}}={}&H_5-\frac{90}{\pi}\mathcal C_\pi(\alpha).
   \label{x_reciprocidad:eq:el-eta}
\end{align}
These are subsequent readings of dodecaphase closure and its regional incidence. The integer \(169\) is retained as the regional selector already constructed, and the coefficients of \(H_5\) as those of its local character; they are not replaced by an interpolation of the Planck constant. The formulas specify exactly which readings are needed for the electronic stage. Other mass-theory functionals need not be introduced here.

For \(0<\alpha<1\), we have \(0<D_A<1\), so \(D_A\alpha<1\). The rational term of \eqref{x_reciprocidad:eq:el-cpi} preserves the complete sum of the geometric tail
\begin{equation}
 \frac{D_A\alpha^7}{1-D_A\alpha}
 =\sum_{n=0}^{\infty}D_A^{n+1}\alpha^{n+7}.
 \label{x_reciprocidad:eq:el-cola-regional}
\end{equation}
It is not a truncation that can be silently modified when precision changes. After the term of index \(N\), the remaining tail is \(D_A^{N+2}\alpha^{N+8}/(1-D_A\alpha)\), allowing evaluation to be controlled without recourse to a target quantity.

\begin{definition}[Action sections and return reader]
The sections before and after return are
\begin{equation}
 \hbar_{\mathrm{pre}}=H_5\,10^{-34}\mathcal U_S,
 \qquad
 \hbar_{\mathrm{ret}}=\eta_{\mathrm{ret}}\,10^{-34}\mathcal U_S.
 \label{x_reciprocidad:eq:el-secciones-accion}
\end{equation}
Whenever \(\eta_{\mathrm{ret}}>0\), define
\begin{equation}
 \delta_{\mathrm{ret}}=H_5-\eta_{\mathrm{ret}},\qquad
 R_{\mathrm{act}}=\frac{\hbar_{\mathrm{pre}}}{\hbar_{\mathrm{ret}}}
 =\frac{H_5}{\eta_{\mathrm{ret}}}.
 \label{x_reciprocidad:eq:el-ratio}
\end{equation}
\end{definition}

\begin{proposition}[Positivity and invariance of the ratio]
In the intervals used in the proof of \eqref{x_reciprocidad:eq:el-decada}, \(0<\eta_{\mathrm{ret}}<H_5\), and therefore \(R_{\mathrm{act}}>1\). The ratio does not change when the basis \(\mathcal U_S\) is replaced by any other common positive basis.
\end{proposition}
\begin{proof}
The terms of \(\mathcal C_\pi\) are positive. For an elementary bound, it suffices to use \(\alpha<0.008\), \(\varphi<1.619\), \(\alpha>0.007\), \(D_A<1\) and \(\pi>3\). The root in \eqref{x_reciprocidad:eq:el-h5} exceeds one, whereas the sum of negative terms is less than \(0.008001\); in particular, \(H_5>0.99\). Likewise,
\[
 0<\frac{90}{\pi}\mathcal C_\pi
 <30\left(169(0.008)^6+
           \frac{(0.008)^7}{1-0.008}\right)<10^{-6}.
\]
Thus \(\eta_{\mathrm{ret}}>0.989999>0\) and \(\eta_{\mathrm{ret}}<H_5\). Finally, the change of basis divides both action coordinates by the same positive factor, which cancels in their ratio.
\end{proof}

Subsequent evaluation gives
\begin{align*}
 H_5&=1.0545718176461544696\ldots,\\
 \eta_{\mathrm{ret}}&=1.0545718169150390611\ldots,\\
 R_{\mathrm{act}}&=1.0000000006932817630\ldots.
\end{align*}
The difference and quotient are two readings of the same sections, one additive and the other multiplicative. They do not represent two independent Planck constants. The passage from angular action to action per complete turn is expressed in each section by \(h=2\pi\hbar\). The subsequent chart \(\mathcal U_S\mapsto\mathrm{J\,s}\) assigns a unit but does not act retroactively on the decimal order already proved.

Regional return can also be written in an angular chart. This reparametrization is not an indispensable additional dependency: \eqref{x_reciprocidad:eq:el-h5}--\eqref{x_reciprocidad:eq:el-eta} supply \(H_5\), \(\eta_{\mathrm{ret}}\) and \(R_{\mathrm{act}}\) directly. If angle and counter-angle are introduced, both must be transported in the same chart; degrees and radians are not mixed within a difference.


% BEGIN OWNER /Users/ruben/Documents/New project/output/EXTREMA_MEDIA_RAZON_RECIPROCIDAD_ES_EN_20260918/ARTICULO/ES/source/sections/revision_planck.tex
\subsubsection{Regional continuation and renewal equation}

The selector \(169\) determines the degree-six impulse. Subsequent extension is organized by the transfer \(D_A\) and unit normalization of the determinant line. These rules admit a formal-series formulation before substituting \(z=\alpha\), with \(D_A\) already constructed. Write
\[
 \mathscr C(z)=169z^6+R(z),\qquad R(z)\in z^7\mathbb R[[z]].
\]
Concatenating an extension increases degree by one and multiplies its weight by \(D_A\). The first strict extension has weight \(D_A\). Its renewal equation is therefore
\begin{equation}
 R(z)=D_Az^7+D_AzR(z).
 \label{x_reciprocidad:eq:revision-renovacion}
\end{equation}

\begin{proposition}[Formal uniqueness of the normalized continuation]
Equation \eqref{x_reciprocidad:eq:revision-renovacion} determines a unique series, and
\[
 \mathscr C(z)=169z^6+\frac{D_Az^7}{1-D_Az}
             =169z^6+\sum_{n=7}^{\infty}D_A^{n-6}z^n.
\]
The realization is analytic in \(|D_Az|<1\).
\end{proposition}
\begin{proof}
The element \(1-D_Az\) has constant term one and is invertible in the formal power-series ring. Solving the equation produces the displayed series. Equivalently, the degree-seven coefficient is \(D_A\) and each subsequent coefficient is \(D_A\) times the preceding one. The geometric series converges absolutely in the stated disk.
\end{proof}

Normalization of the first extension is an indispensable operational datum. If only the degree-six impulse and concatenation ratio were retained, the series
\[
 169z^6+\lambda\frac{D_Az^7}{1-D_Az}
\]
would still share both data for any \(\lambda\). The unit rule fixes \(\lambda=1\) before evaluation. Each rule's contribution to reader uniqueness is thus located, and the rational expression preserves the complete continuation.

\subsubsection{Reduced Planck constant and return of the action section}

In the dimensional chart of \eqref{x_reciprocidad:eq:el-secciones-accion}, the section before return is the model's construction intended for comparison with the reduced Planck constant:
\[
 \hbar_{\mathrm{pre}}=H_5\,10^{-34}\mathcal U_S.
\]
The subsequent section also preserves the regional compensation \(90\mathscr C(\alpha)/\pi\). Both belong to the same action line, and their quotient \(R_{\mathrm{act}}\) measures the multiplicative effect of return. The passage to \(h\) through \(2\pi\) uses the relationship between angular parametrization and a complete turn.

This construction effectively enters \(R_{\mathrm{act}}^{80-54\alpha+6\Delta_4}\). The electron preserves its central coordinate and trajectory; the chronological scale does not replace this structure with an undifferentiated quantum. Nor is a mass already expressed in MeV multiplied again by \(10^{-34}\): the decade belongs to the action section, cancels in its dimensionless quotient and reappears only where the corresponding dimensional realization acts.

The numerical comparison of \(H_5\), the returned section and \(h/(2\pi)\) is presented in Section~\ref{x_reciprocidad:sec:revision-planck-contraste}. This distinguishes the name of the physical object, the equation assigned to it by the model and the precision with which that equation reproduces it.

% END OWNER


\subsection{Two readers of the electronic return}

Electronic memory is represented by a tritic word of length
\(108\) and a signed twelve-phase word. To make the exponent
calculation reproducible, we specify the chronological projection
used. Initially, let
\((r_0,c_0)=(5,5)\) and \((h_0,v_0)=(1,1)\). At each step, the selector
repeats \((+1,-1,0)\). With orientation fixed within each
twenty-seven-step block, the visible rule is
\[
 \begin{array}{c|c|c}
 \text{mode}&\text{read integer}&\text{cellular transport}\\\hline
 +1&r+c&c\mapsto1+(c-1+h)\bmod9\\
 -1&rc&r\mapsto1+(r-1+v)\bmod9\\
 0&9&\text{without displacement}
 \end{array}
\]
The emitted symbol is \(\rho_9(\text{read integer})\bmod3\).
After steps \(27,54,81\), \(h\) and \(v\)
are inverted simultaneously. This rule calculates the observable projection; the
state update additionally preserves residue, quotient,
carry, sheet, boundary and memory. The register of
\(108\) symbols is not identified with the complete state.

In each complete three-step block, each coordinate advances
once; nine such blocks are needed to return to the center in
the same phase. The minimal length of this recorded return is
therefore \(27\). If phase is forgotten, the cell is already reached after
step \(26\), before the neutral threshold completing the block; that is not
the typed return counted by \(\mathcal K_e\). The explicit reading
of the first two return blocks gives
\begin{equation}
 a=100000220,\qquad b=120200000,\qquad
 \gamma_e=(a^3b^3)^2.
 \label{x_reciprocidad:eq:el-palabra}
\end{equation}
Here, powers denote word concatenation. In particular,
\(\gamma_{e,t+54}=\gamma_{e,t}\). In the record of nine steps per
phase, the central orientations give
\begin{equation}
 \tau_e=(1,1,1,-1,-1,-1,1,1,1,-1,-1,-1).
 \label{x_reciprocidad:eq:el-tau}
\end{equation}
The periodicity of these representations does not amount to restarting the
history that produces them.

\begin{definition}[Cyclic discrepancy reader]
For a word \(\gamma\in\mathbb F_3^{108}\), define
\[
 D_k(\gamma)=\#\{t\in\mathbb Z/108\mathbb Z:
                   \gamma_t\ne\gamma_{t+k}\}.
\]
On the domain of words whose period divides \(54\), the integer
reader is
\begin{equation}
 K_D(\gamma)=\left(D_{27}+D_{36},D_1,\frac{D_4-D_3}{2}\right).
 \label{x_reciprocidad:eq:el-kd}
\end{equation}
\end{definition}

\begin{lemma}[Integrality and central evaluation]
The \(D_k\) on the preceding domain are even. For
\eqref{x_reciprocidad:eq:el-palabra},
\[
 \begin{array}{c|rrrrr}
 k&1&3&4&27&36\\\hline
 D_k&54&56&68&48&32
 \end{array}
\]
y \(K_D(\gamma_e)=(80,54,6)\).
\end{lemma}
\begin{proof}
The map \(t\mapsto t+54\) pairs discrepancy indices without fixed
points and preserves their condition; each cardinality is even. For the
explicit word \(a^3b^3\) of length \(54\), the discrepancy
count for each displacement is, in the same order,
\(27,28,34,24,16\). Doubling the word doubles these five cardinalities.
Finally,
\[
 (D_{27}+D_{36},D_1,(D_4-D_3)/2)
 =(48+32,54,(68-56)/2)=(80,54,6).
\]
\end{proof}

The shifts \(27\) and \(36\) are the chronological lifts
of the channels of \(90\) and \(120\) degrees in a chart
where nine steps correspond to thirty degrees. The angular reading
explains that representation, but the integers are computed on the
word, without introducing a mass or a sought real value.

\begin{definition}[Dodecaphasic window reader]
For \(\tau\in\{-1,0,1\}^{12}\), with cyclic indices, let
\begin{align*}
 C_k(\tau)&=\sum_{r=0}^{11}\tau_r\tau_{r+k},\\
 A_\ell(\tau)&=\sum_{r=0}^{11}
           \left(\sum_{j=0}^{\ell-1}\tau_{r+j}\right)^2,\\
 P_6(\tau)&=\sum_{r=0}^{5}\tau_r\tau_{r+6}.
\end{align*}
We define
\begin{equation}
 \Omega(\tau)=4A_3(\tau)-3A_4(\tau),\qquad
 K_\Omega(\tau)=(\Omega(\tau),54,P_6(\tau)).
 \label{x_reciprocidad:eq:el-komega}
\end{equation}
The \(54\) in this reader records the declared chronological half-turn.
It is not confused with the local APP jump of value four.
\end{definition}

\begin{lemma}[Primitive contrast and correlation form]
Among integer linear contrasts \(uA_3+vA_4\) that vanish
whenever \(A_3=3a\) and \(A_4=4a\), the primitive pair, with
orientation \(90-120\), is \((u,v)=(4,-3)\). Moreover,
\begin{equation}
 \Omega=-2C_1-4C_2-6C_3.
 \label{x_reciprocidad:eq:el-omega-corr}
\end{equation}
\end{lemma}
\begin{proof}
Vanishing for every \(a\) requires \(3u+4v=0\). By coprimality,
the solutions are \((u,v)=n(4,-3)\), and orientation fixes the sign
of the primitive generator. Expanding the squares gives
\[
 A_\ell=\ell C_0+2\sum_{k=1}^{\ell-1}(\ell-k)C_k.
\]
Therefore,
\(A_3=3C_0+4C_1+2C_2\) and
\(A_4=4C_0+6C_1+4C_2+2C_3\). Their combination \(4A_3-3A_4\)
cancels \(C_0\) and yields \eqref{x_reciprocidad:eq:el-omega-corr}.
\end{proof}

The uniqueness proved belongs to that class of integer linear
contrasts; it does not by itself exclude every nonlinear functional of the windows.
Specifying the functional class avoids turning an identity
of two channels into a claim of universal selection.

\begin{theorem}[Return triple of the electronic section]
The two readers applied to the records of \(\Gamma_e\)
coincide and give
\begin{equation}
 K_D(\gamma_e)=K_\Omega(\tau_e)=(80,54,6).
 \label{x_reciprocidad:eq:el-triple}
\end{equation}
Their evaluation in the already generated coordinates is
\begin{equation}
 \kappa_e=80-54\alpha+6\Delta_4.
 \label{x_reciprocidad:eq:el-kappa}
\end{equation}
\end{theorem}
\begin{proof}
The first reader was computed above. For \(\tau_e\), the antipodal
product is always one, and hence \(P_6=6\). The cyclic count gives
\(C_1=4\), \(C_2=-4\), \(C_3=-12\); thus
\[
 \Omega=-2\cdot4-4(-4)-6(-12)=80.
\]
Equivalently, \(A_3=44\), \(A_4=32\), and
\(4\cdot44-3\cdot32=80\). Both triples are therefore
\((80,54,6)\). Evaluating the functional
\((A,B,C)\mapsto A-\alpha B+\Delta_4C\) proves the last formula.
\end{proof}

This coincidence concerns the central section. It does not identify
the two readers on every route domain: one compares tritic
symbols and the other compares windows of a signed record. Nor is
\(80\) obtained from the period of a word assigned to Euler's
number. Its origin in this composition is the triple of the electronic
route.

\subsection{Complete electronic composition and normalization}

\begin{definition}[Multiplicative return]
The action of return memory on the basal coordinate is
\begin{equation}
 \varepsilon_e^\beta
 =\varepsilon_e^{(0)}\exp\!\left(\kappa_e\log R_{\mathrm{act}}\right).
 \label{x_reciprocidad:eq:el-beta}
\end{equation}
\end{definition}

\begin{theorem}[Positive electronic factorization]
With the declared readers, orientation, and normalization,
\eqref{x_reciprocidad:eq:el-beta} coincides with \eqref{x_reciprocidad:eq:el-formula-completa}. All
its factors are dimensionless and the resulting coordinate is positive.
\end{theorem}
\begin{proof}
The incompatibility norm is \eqref{x_reciprocidad:eq:el-prefactor}; modal transfer,
deficit, and retained memory compose \eqref{x_reciprocidad:eq:el-basal}. The
triple \eqref{x_reciprocidad:eq:el-triple} produces \eqref{x_reciprocidad:eq:el-kappa}, and the ratio of
the two action sections is \eqref{x_reciprocidad:eq:el-ratio}. Substituting those
three identities into \eqref{x_reciprocidad:eq:el-beta} yields exactly
\eqref{x_reciprocidad:eq:el-formula-completa}. The basal stage is positive, as is the
exponential of a real number. The common factor with the dimension of
action has already canceled in forming \(R_{\mathrm{act}}\).
\end{proof}

The evaluation of this composition begins with
\begin{equation}
 \varepsilon_e^\beta
 =0.5109989506900008086082571648\ldots.
 \label{x_reciprocidad:eq:el-beta-decimal}
\end{equation}
The digits are a subsequent evaluation of the formula, not data used
to select \(22\), \(15\), or \((80,54,6)\). The numerical
representation of the basal coordinate by a finite prefix introduces only its
truncation error; it does not modify the exact formula defining it.

The effect of a different torsional normalization can be located
without comparison to a measurement. Write
\(B=\frac{\sqrt3}{4}[\exp(\varphi/\pi^2)-22\alpha^3]\),
\(R=R_{\mathrm{act}}\), and
\[
 \mathcal E(d)=B(1+15d)R^{80-54\alpha+6d}.
\]
For the two defects in \eqref{x_reciprocidad:eq:el-delta},
\begin{equation}
 \frac{\mathcal E(\Delta_4)}{\mathcal E(d_4)}
 =\frac{1+15\Delta_4}{1+15d_4}
   R^{6(\Delta_4-d_4)}>1.
 \label{x_reciprocidad:eq:el-normalizacion-beta}
\end{equation}
The first torsional inequality and \(R>1\) prove that they are not the same
coordinate. If only the defect in the exponent is changed, keeping the
complete basal coordinate, the ratio of the two outputs is
\(R^{6(\Delta_4-d_4)}\), again different from one. A very small
difference is still a difference of formula; it cannot be attributed to
a mere conversion of units.

\subsection{Dimensional realization: energy, mass, and chronological scale}

The electronic coordinate finally acts on a dimensional line.
Let \(\mathcal L_E\) be an energy line with positive basis
\(\mathcal U_E\). The realization is
\begin{equation}
 E_e^{(0)}=\varepsilon_e^{(0)}\mathcal U_E,\qquad
 E_e^\beta=\varepsilon_e^\beta\mathcal U_E,
 \qquad m_e^\beta=\frac{E_e^\beta}{c_{\mathrm{int}}^2}.
 \label{x_reciprocidad:eq:el-dimensional}
\end{equation}
The last expression belongs to the mass line when
\(c_{\mathrm{int}}\) carries the dimension of speed in the declared
realization. In the comparative chart \(\mathcal U_E=1\,\mathrm{MeV}\),
\eqref{x_reciprocidad:eq:el-beta-decimal} expresses an energy in MeV; the corresponding
mass is expressed in \(\mathrm{MeV}/c^2\). It is common
to abbreviate both quantities by taking \(c=1\), but that abbreviation must not
obscure the dimensional map.

The action scale also admits an action--clock realization. If
\(t_0\) is a positive temporal section, the \(108\)-step clock
defines
\begin{equation}
 \omega_0=\frac{2\pi}{108t_0},\qquad
 E_{\mathrm{clk}}=\hbar_{\mathrm{ret}}\omega_0,\qquad
 m_{\mathrm{clk}}=E_{\mathrm{clk}}c_{\mathrm{int}}^{-2}.
 \label{x_reciprocidad:eq:el-cuanto-reloj}
\end{equation}
Action times frequency has the dimension of energy. This chronological
quantum is a common dimensional section; it is not, by definition,
the electronic section. The route word and its factors are not
obtained from the uniformity of the clock.

\begin{proposition}[Cancellation of the common scale]
The factor \(10^{-34}\mathcal U_S\) in
\eqref{x_reciprocidad:eq:el-secciones-accion} is not part of \(\kappa_e\). On
the same energy line,
\begin{equation}
 \frac{E_e^\beta}{E_e^{(0)}}=R_{\mathrm{act}}^{\kappa_e}.
 \label{x_reciprocidad:eq:el-razon-energia}
\end{equation}
The result is independent of the choice of the common dimensional
basis.
\end{proposition}
\begin{proof}
In forming \(\hbar_{\mathrm{pre}}/\hbar_{\mathrm{ret}}\), the factor
\(10^{-34}\mathcal U_S\) cancels exactly. Dividing the two
energy expressions in \eqref{x_reciprocidad:eq:el-dimensional} also cancels
\(\mathcal U_E\), and \eqref{x_reciprocidad:eq:el-beta} gives the stated ratio.
\end{proof}

Therefore, a coordinate already expressed in MeV is not multiplied again
by \(10^{-34}\). If \(E_{\mathrm{clk}}\) is chosen as the basis of
\(\mathcal L_E\), the coordinates must be transported by the change-of-basis
factor; retaining unchanged the numerical value from the one-MeV chart
would identify two bases without proving it. This is the precise point
at which the action scale enters the absolute realization and ceases to
enter into the relative corrector.

There is also a linear realization that makes the separation
between clock and genealogy visible. On
\(\ell^2(\mathbb Z/108\mathbb Z;\mathbb R)\), the cyclic
shift \(S\) has fixed space
\[
 \ker(S-I)=\mathbb R n_0,\qquad
 n_0=\frac1{\sqrt{108}}\sum_{t=0}^{107}\delta_t.
\]
If \(g_{\Gamma_e}\) is the route symbol in the free linearization
of genealogies and \(\mathcal U_M\) is a basis of the mass line,
the rank-one map is defined by
\begin{equation}
 \mathfrak C_e:
 \mathbb R n_0\otimes\mathbb R g_{\Gamma_e}\longrightarrow\mathcal L_M,
 \qquad
 \mathfrak C_e(n_0\otimes g_{\Gamma_e})
 =\varepsilon_e^\beta\mathcal U_M.
 \label{x_reciprocidad:eq:el-rango-uno}
\end{equation}
The map assembles the preceding factorization once its bases have been declared.
The uniform vector \(n_0\) does not select \(\varepsilon_e^\beta\):
that coordinate comes from APP, the oriented return, and its readers.
An algebraic realization of the null mode is thus distinguished from an
identification of the electron with the clock quantum.

The internal conclusion is the positive factorization of a central
section, with its explicit operations, memory and normalizations.
Assignment to a physical electron adds the charge representation,
spin, coupling and corresponding units chart; it does not
follow exclusively from a decimal coincidence. Metrological
comparison is presented after the exceptional prolongation and does not alter
the inputs to any construction performed here.


% BEGIN OWNER /Users/ruben/Documents/New project/output/EXTREMA_MEDIA_RAZON_RECIPROCIDAD_ES_EN_20260918/ARTICULO/ES/source/sections/revision_electron_argumento.tex
\subsubsection{Structural dependence of the electronic coordinate}

The electronic equation expresses a hierarchical composition. Its first
factor, \(\sqrt3/4\), belongs to the central matrix structure and fixes
a geometric normalization. The exponent \(\varphi/\pi^2\) uses the
already generated regional coordinates in the indicated orientation of the
areal--volumetric reader. The term \(22\alpha^3\) incorporates the dodecaphasic
coupling into the cubic correction. The factor \(1+15\Delta_4\) preserves
the global torsional normalization. Finally, the quotient of the action
sections acts with the exponent determined by the return readers.
Each operation thus has a different antecedent within the same
construction.

The articulation becomes more precise when its dependencies are considered.
The central norm can be calculated before the regional coordinates
are evaluated. The exponential and cubic correction require those
coordinates and the closure of \(\alpha\). Multiplicative memory additionally
requires the action sections and the \(108\)-position trajectory.
Knowing the basal expression therefore does not amount to having evaluated
the complete coordinate. The succession of operations does not add an
external mass to the point \((5,5)\); it develops the readers assigned to its
persistent section.

The distinction between scale and structure also provides an
interpretation of the physical comparison. Mass is the dimensional
realization of the complete coordinate, while the centre, involution,
and vacancy describe the relations that produce it. A common change
of unit modifies the dimensional coordinates, but leaves the dimensionless
ratios and the proved transport identities unchanged.
This is the sense in which the electron preserves its structural character
when expressed relative to a chronological scale.

% END OWNER



% BEGIN OWNER /Users/ruben/Documents/New project/output/EXTREMA_MEDIA_RAZON_RECIPROCIDAD_ES_EN_20260918/ARTICULO/ES/source/sections/realizacion_electromagnetica.tex
\newpage
\subsection{Lorentzian transport and electromagnetic representation}

The electronic section provides a structure and a coordinate
before its application to an interaction. The Lorentzian realization
of cubic incidence determines, with the declared temporal convention,
the space of forms in which that interaction can be expressed.
To formulate it, an oriented local realization \((\mathcal M,g)\)
of dimension four is fixed, with tangent fibers identified with the model
\(Q_\square\). A connection \(\nabla\) with
\(\nabla g=0\), a potential one-form \(A\in\Omega^1(\mathcal M)\),
its curvature \(F=dA\) and a charge representation of the section are also specified.
These data are the arguments of the physical representation that follows.

Hodge duality satisfies \(\star^2=-I\) on two-forms.
The electric and magnetic separation corresponds to a temporal choice
in the decomposition \(1+3\); \(F\) and \(\star F\) preserve their
relation within the same space of two-forms. Neutrality of a
charge representation means \(q=0\), while the section may
preserve orientation, phase and internal transport. Absence of
electric coupling in that representation has content distinct
from absence of phase information.

In this dimensional realization, let \(m_e>0\), \(q\) be the charge,
\(\gamma(s)\) a parametrized trajectory and \(u=\dot\gamma\)
its four-velocity. The Lorentz force is expressed by
\begin{equation}
 m_e\nabla_u u^\mu=qF^\mu{}_{\nu}u^\nu.
 \label{x_reciprocidad:eq:lorentz-electron}
\end{equation}
Antisymmetry of \(F\) implies
\(F_{\mu\nu}u^\mu u^\nu=0\). Consequently,
\[
 \frac{d}{ds}g(u,u)=2g(\nabla_u u,u)=0.
\]
Normalization of the trajectory is preserved provided \(\nabla\)
is metric and \eqref{x_reciprocidad:eq:lorentz-electron} holds. The mass used
is the dimensional realization of the preceding electronic section;
the coupling parameter is identified in the adopted electromagnetic
convention. The specific derivation of the potential and a field
configuration constitutes an additional problem with its own boundary
conditions. The result presented here determines the domain, representation
data and conservation implied by its equation.

TPK route inversion and orientation of an electronic line
then offer a precise comparison: each has its involution,
transport and representation. The Wheeler and Feynman antecedent
places the question historically; equivalence of two realizations
is decided through their maps and compatibility conditions.
This arrangement allows the arithmetic origin of the electron to remain connected
to the formulation of its interaction, making explicit the function
of each stage.

% END OWNER

